\chapter{The Ruler Algorithm}
\label{chap:ruler-algorithm}

% Closed question: What reference structure lets distinct observers compare
% their counts as one physical record?

\begin{chapteressay}

Every measurement assumes a ruler. A measurement of length assumes a
calibrated meter stick. A measurement of mass assumes a calibrated
balance. A measurement of time assumes a calibrated clock. The ruler
is the reference; the measurement is the comparison of the measured
object against the reference.

If the ruler is wrong, the measurement is wrong --- but in a particular
way. Two observers using the same (wrong) ruler can still agree with
each other; they share a calibration error but their comparisons among
themselves remain internally consistent. Two observers using different
(correct) rulers may disagree if the rulers have not been mutually
calibrated. The agreement among observers depends on the shared
calibration, not on the absolute correctness of the calibration.

This is the chapter's question. Two observers, perhaps in different
locations, perhaps using different physical instruments, perhaps
working at different times, want to compare their measurements. What
structure must they share to make the comparison admissible?

The answer is a calibrated reference. The reference is the ruler. Both
observers must measure against the same ruler. If they cannot agree on
the same physical ruler (perhaps because they are in different
locations and cannot ship the ruler back and forth), they must agree on
a calibration chain: a sequence of comparisons that links each
observer's local ruler to a common reference.

The chapter develops this through five algorithmic forms. The first is
tail-latency measurement: a web service measures its 99th percentile
latency and uses it as the ruler for SLA compliance. The percentile is
the calibrated reference; new request latencies are compared against
it. The second is Dijkstra's algorithm: the shortest path in a network
is computed by comparing edge weights against a calibrated distance
function; the function is the ruler. The third is the Lean
\texttt{EKG} typeclass: the calibration is a private \texttt{Reading}
structure exposed only through propositional comparisons. The fourth
is Gram-Schmidt orthogonalization: building a calibrated orthonormal
basis one vector at a time. The fifth is the coordinate transform:
converting measurements between two calibrations of the same physical
quantity.

The second concrete example is the factory quality-control room. A
metrology technician uses a set of precision gauge blocks --- pieces of
hardened steel ground to specific thicknesses --- as the factory's
reference. Every measurement made in the QC room compares the part
being measured against the gauge blocks. The gauge blocks themselves
are calibrated against a secondary reference, which is calibrated
against the national standard, which is calibrated against the
international Bureau of Weights and Measures. The chain anchors at
the speed of light in vacuum and a specific time interval.

The QC room is the practical embodiment of the chapter's structure.
The calibration is hidden inside the instruments (the micrometer's
gear train, the gauge blocks' precise dimensions, the granite
surface table's flatness). The interface is the propositional
comparison (in tolerance / out of tolerance). The calibration chain
reaches up to a universal reference.

The same structure appears in many domains. In software, the
calibration is the version control system that ensures two
developers see the same source code. In finance, the calibration is
the currency exchange rate that converts measurements in dollars to
measurements in euros. In medicine, the calibration is the reference
range that determines whether a blood test is normal or abnormal.
In each case, the algorithm depends on the shared reference;
measurements are admissible only through the reference.

The chapter's closed question is:

\begin{center}
\emph{What reference structure lets distinct observers compare their
counts as one physical record?}
\end{center}

The answer the chapter develops is the ruler algorithm with a
hidden calibration. The calibration is a reference scale that is
not directly exposed to observers; observers interact with it only
through propositional comparisons. The hidden-ness is intentional:
direct access to the scale would tempt observers to make non-
comparison-based judgments, which the algorithm cannot certify.
The propositional interface forces all comparisons through the
shared, calibrated reference.

The Lean \texttt{EKG} type is the chapter's clean formal
representation. The \texttt{Reading} structure is private. The
\texttt{boolean} comparison is the public interface. The
\texttt{LOGICAL} typeclass requires the calibration; the
\texttt{ComputerProgram.le} comparison routes through it. The
compiler enforces the protocol: programs that bypass the
calibration cannot be admitted.

The chapter's claim is that this is the algorithmic form forced by
the chapter's question. Any system that compares distinct
observers' measurements must have a shared, calibrated reference.
The reference is hidden; the comparisons are propositional. The
discipline is what makes the comparisons admissible.

Subsequent chapters build on this. Chapter 9 addresses what
happens when two transports through the calibrated frame do not
commute: the curvature, the commutation failure, the algorithmic
record of paths that disagree. Chapter 10 addresses what survives
admissible recompilation: the invariants of the algorithmic
description. Chapter 11 closes with the append-only nature of the
final ledger and the type-check that ends the book.

\end{chapteressay}

\section{Tail-Latency Measurement}
\label{se:tail-latency-measurement}

A web service processes ten thousand requests per second. Each request
is timed: the server records, in microseconds, how long the request
took to handle. At the end of each minute, the server has six hundred
thousand latency measurements. The question is: what is the latency
of this service?

The naive answer is the mean: average all six hundred thousand
measurements and report the result. But the mean hides the worst
cases. If 99 percent of requests are fast (under ten milliseconds)
and 1 percent are very slow (several seconds), the mean might be a
few hundred milliseconds, which understates how bad the bad cases
are. The mean is misleading.

The standard answer is percentile-based: sort all measurements,
report several key percentiles. The 50th percentile (median) is the
middle latency. The 99th percentile is the latency below which 99
percent of requests fall. The 99.9th percentile is the latency
below which 99.9 percent fall. The 99.99th percentile is the latency
of the worst out of ten thousand.

For SLA (service-level agreement) purposes, the 99th percentile is
often the contractual reference. The agreement says: the service
will respond in under 500 milliseconds for 99 percent of requests.
The measurement instrument is the percentile-based latency analysis;
the reference is the SLA threshold. Compliance is a comparison
against the reference.

This is the chapter's question in operational form. The 99th
percentile is the ruler: the calibrated reference against which
individual requests are judged. A request whose latency is below
the ruler value is compliant. A request whose latency is above is
non-compliant. The ruler is set by the SLA, not by the request's
own characteristics. The ruler is the system-wide calibration.

Computing the percentile requires sorting the latencies, which is
$O(n \log n)$ for $n$ measurements. For six hundred thousand
measurements, this is about $10^7$ operations: a small fraction of
a second on a modern CPU. The sorting is the algorithm; the
percentile is the algorithm's output. The output is the ruler value.

The ruler can be applied to new measurements as they arrive. Each
new request is timed; the timing is compared to the ruler; the
request is either compliant or not. The comparison is
$O(1)$ per request; the percentile is recomputed periodically (every
minute, every hour, every day, depending on the system's
requirements).

The Lean analog is the comparison operation routed through the
\texttt{EKG} typeclass. Each program's elaboration cost is recorded
as a heartbeat count. The \texttt{EKG} comparison asks: is this
program's cost within tolerance of the reference cost? The
reference cost is the calibration --- the ruler. The comparison is
the application of the ruler to the program.

The \texttt{UniverseTensor} type in \texttt{Episode10.lean} extends
the structure to multi-dimensional reference. A
\texttt{UniverseTensor} represents the calibrated reference for an
entire system: not just one ruler, but a tensor of rulers, one for
each dimension being measured. The system's calibration is the
tensor; comparisons in each dimension route through the corresponding
ruler.

The chapter's claim is that the ruler algorithm is the algorithmic
form forced by the chapter's question. To compare distinct
observers' counts as one physical record, the system must have a
shared reference. The reference is the ruler; the comparison is the
ruler's application. The ruler is calibrated; the application is
deterministic. Different observers using the same ruler produce
comparable measurements. Different rulers (different SLA percentiles,
different calibration standards) produce incomparable measurements
that the system cannot unify without a recalibration step.

\section{Dijkstra as Ruler Algorithm}
\label{se:dijkstra-as-ruler-algorithm}

Dijkstra's algorithm finds the shortest path from a source node to all
other nodes in a weighted graph with non-negative edge weights. The
algorithm maintains a priority queue of nodes ordered by their current
known minimum distance from the source. At each step, it extracts the
minimum-distance node from the queue, marks it as visited, and updates
the distances of its unvisited neighbors based on the new
shortest paths through the visited node.

For a network of one thousand servers with five thousand
connections, Dijkstra finds the shortest path from any source in
$O((V + E) \log V) \approx 8 \times 10^4$ operations: a small
fraction of a second.

The algorithm's correctness depends on the ruler. The ruler is the
distance function: it assigns to each edge a non-negative weight,
and the path distance is the sum of edge weights along the path.
The algorithm always extracts the smallest-distance node next.
This is admissible because of the ruler's properties: distance is
non-negative, additive, and gives a total order on partial paths.

Without the ruler, the algorithm cannot make a deterministic
choice about which node to extract. With it, the choice is
forced: the smallest-distance node is unambiguous. The ruler is
the algorithm's calibration. Different rulers (different edge
weight functions) produce different shortest paths. The same
graph, with the same algorithm, can produce different answers if
the ruler is different.

This is the chapter's structure in the network domain. The
calibration is the ruler (the edge weight function). The
comparison is the algorithm's extract-min operation. The
admissible movement is along the shortest path. Different
calibrations produce different admissible movements.

A practical example: routing internet packets. The router knows
the network topology and the link delays. Dijkstra computes the
shortest path (in milliseconds of delay) from the router to each
destination. The shortest path is the route. The calibration
(the delay function) is what makes the route admissible.

Changing the calibration changes the routing. If link delays
spike on a particular path, the calibration changes, and Dijkstra
re-runs with the new weights. The new shortest path may avoid the
congested link entirely. The recalibration produces new admissible
routes.

In Lean, Dijkstra's algorithm appears implicitly in the
\texttt{Truth.le} relation in \texttt{Episode10.lean}. The
relation defines a partial order on truths; the shortest-path
computation in this order is the chain of inferences from
hypothesis to conclusion. Each step of the chain has a cost (the
heartbeat count of the inference rule); the shortest chain is
the minimum-cost proof.

The compiler's tactic mechanism, when it searches for a proof,
behaves like Dijkstra: it explores the partial order of truths,
following the lowest-cost edges first, until it reaches the
target. The cost is the proof effort; the calibration is the
\texttt{EKG} reference; the algorithm is the
\texttt{Truth.le}-routed search.

\section{EKG as Calibration Reference}
\label{se:ekg-as-calibration-reference}

The Lean calibration module is in
\texttt{device/Measurement/Calibration/LeanCalibration.lean}. It
defines the \texttt{EKG} type, which represents a calibration's
public interface, and the private \texttt{Reading} structure, which
holds the hidden reference scale.

The \texttt{Reading} structure has a single field:

\begin{leanexcerpt}{Reading}
private structure Reading where
  heartbeats : Nat
\end{leanexcerpt}

The \texttt{Reading} is private. User code cannot directly access
the \texttt{heartbeats} field. The field is the hidden
implementation: the actual count of elaboration heartbeats that
the compiler has performed. The count is a \texttt{Nat}, but the
\texttt{Nat} does not leak into user-facing code.

The \texttt{EKG} structure exposes only \texttt{Prop}-valued
operations:

\begin{leanexcerpt}{EKG}
structure EKG where
  private reference : Reading
  embigger? : Prop $\to$ Prop $\to$ Prop
\end{leanexcerpt}

The \texttt{reference} field is itself private: it holds an
instance of the hidden \texttt{Reading} at the calibration's
baseline. The \texttt{embigger?} method takes two propositions
and returns a proposition asserting that the first dominates the
second under the calibration's bookkeeping. Together they are
sufficient for the chapter's comparison: an \texttt{EKG} can
answer ``is this program at most as costly as that program?''
without exposing either the underlying heartbeat counts or the
baseline.

This is the chapter's central design pattern. The calibration
exists as a hidden reference scale (the \texttt{Reading}, plus
the \texttt{reference} held inside the \texttt{EKG}). The public
interface is propositional comparison (the \texttt{embigger?}
relation). User code interacts with the calibration through the
propositional interface, never with the raw scale. The compiler
enforces the privacy: the \texttt{private} keyword prevents user
code from reading either the heartbeat count or the baseline.

This is the chapter's algorithmic formalization of the metrology
principle that calibration is a reference standard, not a
measurement. The standard exists; comparisons against the standard
are admissible measurements; the standard itself is not directly
exposed. The standard is hidden because exposing it would tempt
users to make non-comparison-based decisions, which the algorithm
cannot certify.

The \texttt{LOGICAL} typeclass carries the \texttt{EKG}:

\begin{leanexcerpt}{LOGICAL}
class LOGICAL
    (Value : Type)
    (Carrier : CarrierProcess Value)
    [DISTINGUISHABLE Value Carrier] ... [UNIVERSAL Value Carrier]
  where
  feelings : HeartbeatProcess Value Carrier
  ekg      : Calibration.EKG
  logical? : YarnTheory $\to$ YarnTheory $\to$ Prop := fun a b =>
    YarnTheory.le a b
\end{leanexcerpt}

A \texttt{LOGICAL} instance is a program that carries its own
calibration. The instance is required for any comparison operation
on \texttt{ComputerProgram} values. Without a \texttt{LOGICAL}
instance, the compiler refuses to admit the comparison: the program
has no calibrated reference, so the comparison is inadmissible.

The chapter's algorithmic claim is that \texttt{EKG} and
\texttt{LOGICAL} together formalize the calibration protocol.
\texttt{EKG} is the reference scale; \texttt{LOGICAL} is the
program's certificate that it carries the scale; comparisons are
routed through the scale via the \texttt{boolean} interface. The
discipline is enforced at compile time.

Without the discipline, programs could compare each other using
ad-hoc rulers --- one program's CPU time against another's memory
usage, with no shared reference. The comparisons would be
incommensurable. The chapter's algorithmic answer is to require
the shared reference (the \texttt{EKG}) and to hide its
implementation (the \texttt{Reading}). The public interface
forces all comparisons through the shared, calibrated reference.

\section{Gram-Schmidt as Sequential Ruler}
\label{se:gram-schmidt-as-sequential-ruler}

Gram-Schmidt appeared in Chapter 7 as an inversion algorithm. In this
chapter it serves a different role: a sequential calibration of an
$n$-dimensional vector space.

The procedure builds an orthonormal basis one vector at a time. The
first basis vector defines the first axis; the second defines the
second axis (orthogonal to the first); the third defines the third
axis (orthogonal to the first two). Each basis vector is a calibration
of one dimension. The full set of basis vectors is the calibration of
the entire space.

Before adding a new basis vector, the algorithm verifies its
calibration: the new vector must be orthogonal to all the previous
basis vectors (its inner products with them are zero). The
orthogonality is the admissibility condition. The verification is the
calibration check.

If the inner product is nonzero, the new vector is not properly
calibrated: its component in the previous basis vectors is removed
by subtraction, and only the orthogonal remainder is admitted as the
new basis vector. The subtraction is the recalibration.

This is the chapter's structure in pure linear algebra. The
calibration is the orthonormal basis. The admissibility check is the
inner product. The recalibration is the subtraction of the
non-orthogonal component. The output is a calibrated basis that
defines admissible measurements in each dimension.

Gram-Schmidt is the calibration-construction algorithm. It does not
measure anything in the input vectors; it constructs the reference
frame within which measurements will be made. The inputs are the
raw data (some independent vectors); the output is the calibrated
basis. Subsequent measurements in the calibrated basis are
admissible because the basis is orthonormal.

The connection to the chapter's question is direct. Two observers
using different bases (calibrations) cannot directly compare their
measurements: their numbers are in different coordinate systems. To
compare, the two bases must be related by a coordinate transform
(the next section). Without the transform, the measurements are
incommensurable.

Gram-Schmidt builds a canonical basis from a starting set of
vectors. The canonical basis is the unique orthonormal basis whose
$k$-th vector lies in the span of the first $k$ input vectors and
has positive inner product with the $k$-th input vector. The
canonical-ness is the calibration's universality: any starting set
that spans the same subspace produces the same Gram-Schmidt
output.

In Lean, the analog is the typeclass cascade
\texttt{Truth} $\to$ \texttt{Variation} $\to$ \texttt{LOCAL}. A
\texttt{Truth} is a propositional statement; \texttt{Variation}
records its first-order sensitivity (its derivative in the
informational sense); \texttt{LOCAL} certifies that the
sensitivity is admissible in the local frame. The cascade builds
the calibration: each step is one Gram-Schmidt step in the
truth-space.

\section{Coordinate Transform as Recalibration}
\label{se:coordinate-transform-as-recalibration}

A thermometer reads 68 degrees. The question is: 68 degrees on
which scale? Fahrenheit reads 68 F at room temperature. Celsius
reads 20 C at the same temperature. Kelvin reads 293 K. The same
physical temperature corresponds to three different numbers in three
different calibrations. To compare a reading from one calibration
to a reading from another, the algorithm must apply a coordinate
transform.

The transform from Fahrenheit to Celsius is:
\[
T_C = \frac{5}{9}(T_F - 32).
\]
The transform from Celsius to Kelvin is $T_K = T_C + 273.15$. The
two transforms compose: $T_K = \frac{5}{9}(T_F - 32) + 273.15$.

The transform is the recalibration. It converts a number in one
calibration to a number in another, preserving the underlying
physical quantity. The same temperature, expressed in any of the
three scales, gives the same physical result (the same molecular
kinetic energy, the same equilibrium with a reference).

The chapter's question is answered here in the simplest possible
form. Two observers, one with a Fahrenheit thermometer and one
with a Celsius thermometer, can compare their measurements by
applying the coordinate transform. The transform is the algebraic
relation between their calibrations. The transform is a single
linear function with a multiplier ($5/9$) and an offset ($-32$).

The transform must be carried with the data. If a measurement is
reported as ``68 degrees,'' the report is incomplete without the
calibration. ``68 F'' and ``68 C'' are different temperatures by
about 20 degrees C. The unit (the calibration) is part of the
measurement, not external to it.

In the Lean codebase, the calibration is preserved through the
\texttt{EKG} typeclass. Programs that use different
\texttt{EKG} instances are not directly comparable; the comparison
must route through a coordinate transform between the two
\texttt{EKG}s. The transform is the algorithmic recalibration.

The general principle: a measurement is a triple (number, unit,
operational context). Two measurements can be compared only after
they are transformed to the same calibration. The transform is
admissible only if its inputs are well-calibrated and its output is
correctly stated.

This is also the principle behind unit conversion libraries.
Software libraries like \texttt{units} in Python or
\texttt{Mathlib.Tactic.NormNum} in Lean enforce dimensional
consistency: an addition of meters and seconds is rejected
because the units are different. A multiplication of meters by
seconds produces meter-seconds, a new unit. The library is the
algorithmic encoder of the recalibration discipline.

The coordinate transform also appears in more sophisticated
contexts. The Lorentz transform in special relativity converts
between two inertial reference frames. The transform preserves
the spacetime interval $c^2 t^2 - x^2$, which is the same physical
quantity in both frames. The transform is not linear in the
ordinary sense; it is a hyperbolic rotation, but the structure is
the same: a coordinate transform that converts between two
calibrations of the same underlying physics.

The chapter's algorithmic message: every transport that crosses a
calibration boundary requires a coordinate transform. The transform
is part of the algorithm. The result is the same physical quantity
expressed in a new calibration. The new calibration's
\texttt{EKG} certifies the result; the transform's correctness
certifies the conversion.

\begin{bridge}

The chapter's question was what reference structure lets distinct
observers compare their counts as one physical record.

The answer is the ruler algorithm. The reference is a calibration:
a hidden scale (the \texttt{Reading} structure in Lean) accessed
only through propositional comparisons (the \texttt{boolean}
operation of \texttt{EKG}). Tail-latency percentiles, Dijkstra's
shortest-path costs, Gram-Schmidt's orthonormal bases, and
coordinate transforms all express the same algorithmic structure:
calibrated rulers that mediate admissible comparisons. The
\texttt{LOGICAL} typeclass certifies that a program carries its
calibration; the \texttt{ComputerProgram.le} comparison routes
through the calibration.

What remains is the failure of transport to commute. Two
calibrated observers carrying the same record through different
paths may produce different results. The difference is the
curvature of the algorithmic transport. Chapter 9 takes up
commutation failure: the algorithmic record of paths that do
not commute.

\end{bridge}

\begin{coda}{The Factory Quality-Control Room}

The quality-control room sits at the end of the production line.
A metrology technician works beneath an array of measuring
instruments, with the factory's reference standards locked in a
cabinet beside her.

The gauge blocks are the heart of the room. Each block is a
hardened-steel piece ground to a specific thickness with a tolerance
of fifty nanometers. By stacking blocks (a technique called
wringing), the technician can produce any thickness in the
factory's working range to within a few hundred nanometers.

The blocks are the factory's calibrated reference. They define what
``1.000 mm'' means inside this building. Every measurement made in
the QC room is compared against the blocks. A part that should be
5.000 mm thick is wrung against the 5.000 mm stack; the difference
is the part's deviation from the standard. The deviation is the
measurement.

This morning's task is the inspection of a batch of three hundred
machined parts. The parts are specified to be 12.450 mm long, with a
tolerance of $\pm 0.010$ mm. The QC technician's job is to verify each
part against the specification.

She begins by checking the calibration of her measuring instruments.
The micrometer is checked against the 5.000 mm gauge block; the reading
is 5.0001 mm, within calibration tolerance. The ring gauge is checked
against the matching plug gauge; the fit is correct. The master square
is checked for perpendicularity against the granite surface table; the
deviation is below the table's certified flatness. All four instruments
pass calibration.

She then measures each of the three hundred parts.

For each part, she places it on the surface table, holds it square, and
brings the micrometer's spindle down to contact the top surface. The
micrometer's reading is a number with five significant figures: 12.451,
12.450, 12.452, 12.453, 12.449. She records each reading on a
clipboard.

The readings are routed through the calibration. The micrometer's
internal mechanism (the spindle, the thimble, the ratchet) translates
the physical contact into a numeric reading. The number is in
millimeters because the micrometer's calibration is in millimeters. The
calibration is hidden: the technician does not see the gear ratios or
the thread pitches; she sees only the displayed number.

This is the chapter's \texttt{EKG} structure in metrology form. The
calibration is hidden inside the instrument. The interface is the
displayed reading. The technician applies the reading; the
\texttt{boolean} comparison is the in-tolerance / out-of-tolerance
judgment.

For each part, the reading is compared to the specification: 12.450
$\pm$ 0.010 mm. A part reading between 12.440 and 12.460 is in
tolerance; a part reading outside this range is out. The comparison
is a single Boolean.

Of the three hundred parts, two hundred ninety-six are in tolerance.
Four are out. Each of the four out-of-tolerance parts is set aside in
a reject bin. The QC technician notes the readings: three parts read
12.439 (just below tolerance, likely a tool-wear issue) and one part
reads 12.461 (just above tolerance, likely a different cause).

The QC report compiles the readings: 296 in tolerance, 4 out, a
specific distribution of within-tolerance readings. The report is
submitted to the production supervisor. The four rejects are returned
to the production floor for analysis: which machine produced them,
when, what was happening at that time.

Toward mid-morning, a new batch arrives from a different supplier.
The supplier's parts are also specified to be 12.450 mm long, but the
supplier uses inch-based calibration internally and converts to
millimeters for the customer. The technician measures the supplier's
parts on the same micrometer and gets readings that cluster around
12.453 mm. All in tolerance, but skewed slightly high.

The skew is the coordinate transform's residue. The supplier's
calibration converts 0.4901 inches to 12.4485 mm via the conversion
factor 25.4 mm/inch. The conversion is exact in principle but
introduces a small bias in practice because the supplier's micrometers
have been calibrated to a slightly different inch standard (the
international inch is 25.4 mm, but historical inch standards differ
by a few parts per million). The bias is small, well within the
$\pm 0.010$ mm tolerance, but it is visible as a systematic offset.

The QC technician notes the offset in her log. The supplier's parts
are accepted, but the systematic offset suggests that the supplier's
calibration chain has a slightly different reference. If the offset
grows, or if the tolerance tightens, the supplier will need to
recalibrate against the factory's reference standards directly. For
now, the offset is within tolerance.

This is the chapter's recalibration discipline in operational form.
Two suppliers use slightly different calibrations of the same physical
quantity (length in millimeters). The coordinate transform between
their calibrations is the conversion factor. The factor is known;
the small bias is the residue of the transform's imperfection. The
QC system tracks the residue and ensures that it stays within
admissible bounds.

In the afternoon, the technician performs a routine recalibration of
the master gauge blocks themselves. She compares them against the
factory's secondary reference: a set of gauge blocks certified by the
national metrology institute. The comparison uses a more precise
instrument, a laser interferometer, which can detect differences down
to a few nanometers. The factory's master blocks are found to agree
with the secondary reference to within 30 nanometers. The agreement
is recorded; the master blocks are certified for another year.

The secondary reference, in turn, is certified by the national
metrology institute, which is certified by the international Bureau
of Weights and Measures. The calibration chain reaches up four
levels: working instrument $\to$ master gauge blocks $\to$
secondary reference $\to$ national standard $\to$ international
standard. At the top, the international standard for the meter is
defined by the speed of light in vacuum and a specific time
interval. The chain anchors at a physical constant.

The chain is the chapter's algorithmic universalization. Every
measurement in the factory is comparable to every measurement in
any other factory that participates in the same chain, because every
measurement traces back to the same international standard. The
universal reference is what makes global manufacturing possible. The
chapter's algorithmic claim is that this universality is precisely
the \texttt{EKG} pattern at scale: a hidden reference (the
international meter), exposed only through calibrated comparisons
(the chain of gauge blocks).

The technician closes the QC room at five o'clock. The day's
measurements are filed. The reject bins are emptied to engineering
for analysis. The calibration log is updated. Tomorrow she will
return to the same instruments, the same gauge blocks, the same
calibration chain. The factory's quality discipline is the
\texttt{LOGICAL} typeclass enforced through metrology: every
measurement carries the calibration; every calibration traces back
through the chain; the chain reaches the universal reference.

\end{coda}
